\documentclass[12pt]{article}

\usepackage[paperwidth=200mm,paperheight=112.5mm,margin=11mm,top=10mm]{geometry}
\usepackage{amsmath,amssymb,booktabs,array}
\usepackage[T1]{fontenc}
\usepackage{lmodern}

\pagestyle{empty}
\setlength{\parindent}{0pt}
\setlength{\abovedisplayskip}{9pt}
\setlength{\belowdisplayskip}{9pt}

\newcommand{\R}{\mathbb{R}}
\newcommand{\T}{\mathsf{T}}
\DeclareMathOperator{\col}{col}
\newcommand{\fbase}{f_{\mathrm{base}}}
\newcommand{\faug}{f_{\mathrm{aug}}}
\newcommand{\Faug}{F_{\mathrm{aug}}}
\newcommand{\ANN}{\mathrm{ANN}}
\newcommand{\panel}[1]{\subsection*{\Large #1}\vspace{1mm}}

\begin{document}

%======================================================================
\panel{Baseline in state-space form}

Physical state $x_b=\col(q,\dot q)\in\R^{6}$ with $q=\col(X,\Theta,Y)$. The
scheduling signal is the $Y$ position itself, $Y=q_3$, so the baseline is
self-scheduled. Written as a first-order system,
$M(Y;v)\ddot q+C(v)\dot q+K(v)q=Pu$ becomes

\begin{equation*}
\dot x_b
=
\left[\begin{array}{cc}
0 & I\\[2.5mm]
-M(Y;v)^{-1}K(v) & -M(Y;v)^{-1}C(v)
\end{array}\right]
x_b
+
\left[\begin{array}{c}
0\\[2.5mm] M(Y;v)^{-1}P
\end{array}\right]
u .
\end{equation*}

One sample of RK4 gives the discrete transition, and only positions are
measured:

\begin{equation*}
x_{b,k+1}=\fbase(v,x_{b,k},u_k),
\qquad
y_k=C_y x_{b,k}+D_y u_k,
\qquad
C_y=\begin{bmatrix}P^{\T} & 0_{3\times 3}\end{bmatrix},
\qquad
D_y=0 .
\end{equation*}

The free parameters are the ten identifiable combinations

\begin{equation*}
v=\col\!\left(k_{b,\Sigma},\,c_{g1},\,c_{g2},\,c_y,\,c_{b,\Sigma},\,m_h,\,
m_{\mathrm{tot}},\,m_{\Delta},\,J_{\mathrm{eff}},\,d\right)\in\R^{10},
\end{equation*}

where sums appear because $k_{b1},k_{b2}$ and $c_{b1},c_{b2}$ and $J_b,J_h$ are
not separately identifiable from the data.

\newpage
%======================================================================
\panel{Additive augmentation in the state equation}

Physical states $x_{b,k}\in\R^{6}$ and learned states $x_{a,k}\in\R^{8}$ are
stacked into $\widehat x_k=\col(x_{b,k},x_{a,k})\in\R^{14}$:

\begin{equation*}
\begin{bmatrix}x_{b,k+1}\\[1.5mm] x_{a,k+1}\end{bmatrix}
=
\begin{bmatrix}\fbase(v,x_{b,k},u_k)\\[1.5mm] 0\end{bmatrix}
+
\begin{bmatrix}
\faug^{b}(\theta_a,\widehat x_k,u_k)\\[1.5mm]
\faug^{a}(\theta_a,\widehat x_k,u_k)
\end{bmatrix},
\qquad\qquad
y_k=C_y x_{b,k}+D_y u_k .
\end{equation*}

One network provides both contributions, split by row:

\begin{equation*}
\Faug(\theta_a,\widehat x_k,u_k)
=\begin{bmatrix}
\faug^{b}(\theta_a,\widehat x_k,u_k)\\[1.5mm]
\faug^{a}(\theta_a,\widehat x_k,u_k)
\end{bmatrix}\in\R^{14},
\qquad
\faug^{b}\in\R^{6},
\qquad
\faug^{a}\in\R^{8}.
\end{equation*}

\newpage
%======================================================================
\panel{Where the two components are added}

Gy\"or\"ok adds them in the output equation,

\begin{equation*}
\widehat y_k=\phi(x_k)\theta_b+f^{\ANN}_{\theta_a}(x_k),
\tag{Eq.~3}
\end{equation*}

where $x_k$ is the lagged input--output regressor, $\phi(x_k)$ the baseline
regressor matrix, $\theta_b$ the physical parameters entering it linearly, and
$\theta_a$ the network weights.

\vspace{2mm}
Here they are added in the state equation, so neither component appears in the
output equation:

\begin{equation*}
\left.
\begin{array}{r}
\fbase(v,x_{b,k},u_k)\\[1.5mm]
\faug^{b}(\theta_a,\widehat x_k,u_k)
\end{array}
\right\}
\ \longrightarrow\ x_{b,k+1}
\ \longrightarrow\ y_{k+1},
\qquad\qquad
\fbase,\ \faug^{b}\ \nrightarrow\ y_k .
\end{equation*}


\end{document}
